% C-29 -- figure merges (figure_merge_audit.md, items A + B-partial + C). Audit copy only --
% live files are sections/methodology.tex and sections/results.tex.

% --- A: methodology.tex, Terreno_irre + Terreno_incli ---
\begin{figure}[t]
  \centering
  \begin{subfigure}[t]{0.48\textwidth}
    \centering
    \includegraphics[width=\linewidth]{TERRENO_IRREGULAR.jpeg}
    \caption{Irregular terrain.}
    \label{fig:Terreno_irre}
  \end{subfigure}\hfill
  \begin{subfigure}[t]{0.48\textwidth}
    \centering
    \includegraphics[width=\linewidth]{TERRENO_INCLINADO.jpeg}
    \caption{Tilted terrain.}
    \label{fig:Terreno_incli}
  \end{subfigure}
  \caption{\revblue{Physical test terrains: (a) irregular Perlin-noise terrain; (b) tilted-incline setup.}}
  \label{fig:terrenos_fisicos}
\end{figure}

% --- C: results.tex, RES_IMU + RES_IMU2 (placed at RES_IMU2's original location) ---
\begin{figure}[t]
  \centering
  \begin{subfigure}[t]{\textwidth}
    \centering
    \includegraphics[width=0.85\linewidth]{RES_IMU.png}
    \caption{Trajectory in an unstable terrain in a simulation trial. The poses were taken every 22 seconds.}
    \label{fig:RES_IMU}
  \end{subfigure}

  \vspace{0.6em}

  \begin{subfigure}[t]{\textwidth}
    \centering
    \includegraphics[width=0.85\linewidth]{RES_IMU2.0.png}
    \caption{Trajectory during traversal of elevated terrain in simulation. Snapshots taken every 12~s.}
    \label{fig:RES_IMU2}
  \end{subfigure}
  \caption{\revblue{Simulated-trajectory snapshot sequences for the two variable-topography scenarios.}}
  \label{fig:res_imu_pair}
\end{figure}

% --- B (partial): results.tex, GRAPH_IMU + GRAPH_IMU2 (placed at GRAPH_IMU2's original location) ---
\begin{figure}[t]
  \centering
  \begin{subfigure}[t]{0.48\textwidth}
    \centering
    \includegraphics[width=\linewidth]{GRAPH_IMU.png}
    \caption{Standard deviation of the acceleration in the Z axis during transport on an unstable terrain.}
    \label{fig:GRAPH_IMU}
  \end{subfigure}\hfill
  \begin{subfigure}[t]{0.48\textwidth}
    \centering
    \includegraphics[width=\linewidth]{GRAPH_IMU2.png}
    \caption{Pitch angle during motion in a steeply sloping terrain.}
    \label{fig:GRAPH_IMU2}
  \end{subfigure}
  \caption{\revblue{IMU signal traces motivating the mode transition in each simulated terrain scenario: (a) unstable/irregular terrain; (b) steeply sloping terrain.}}
  \label{fig:graph_imu_pair}
\end{figure}

% --- B (partial): results.tex, GRAPH_IMU_real + GRAPH_IMU2_real (placed at GRAPH_IMU2_real's original location) ---
\begin{figure}[t]
  \centering
  \begin{subfigure}[t]{0.48\textwidth}
    \centering
    \includegraphics[width=\linewidth]{GRAPH_IMU_real.png}
    \caption{MLP outputs and IMU state timeline on uneven terrain.}
    \label{fig:GRAPH_IMU_real}
  \end{subfigure}\hfill
  \begin{subfigure}[t]{0.48\textwidth}
    \centering
    \includegraphics[width=\linewidth]{GRAPH_IMU2_real.png}
    \caption{MLP outputs and IMU state timeline on an inclined terrain.}
    \label{fig:GRAPH_IMU2_real}
  \end{subfigure}
  \caption{\revblue{MLP-output and IMU-state timelines motivating the mode transition in each physical terrain scenario: (a) uneven terrain; (b) inclined terrain.}}
  \label{fig:graph_imu_real_pair}
\end{figure}

% NOT applied: NEU_IMU/NEU_IMU2 and NEU_IMU_real/NEU_IMU2_real -- excluded, see PROGRESS.md C-29 row
% (their captions' "a"/"b" letters are baked into the PNG itself, not LaTeX subfigure labels).
